\section{Reference implementation, examples, and reproduction}
\label{app:artifacts}

\subsection{What is implemented}

The standard-library-only Python file \code{code/richardson.py} constructs
the expression AST from a supplied integer polynomial. It implements
exact rational interval arithmetic, bounded arctangent and sine series,
outward dyadic rounding, and an a priori precision calculation. It does
not implement MRDP or provide a universal halting polynomial.

The representation shares repeated subexpressions in a directed acyclic
graph. Expanding that graph produces an expression in the grammar of
Section~\ref{sec:statement}. The coefficients $C,D$ and the formulas for
$R,L,F,h_j$ and $E_P$ agree with the mathematical construction.

To certify one strict nonzero value, the program evaluates a rational
interval enclosing $A_P(q)$. The identity detector has the exact monotone
image
\[
 a\longmapsto |a|-a=\max(0,-2a).
\]
Applying that image to the enclosing interval preserves the dependence
between the two occurrences of $A_P$ and gives a tighter enclosure than
independent interval subtraction. Both forms of enclosure are recorded
in the demonstration JSON. A strictly positive lower endpoint certifies
nonidentity. A zero interval at a single point does not certify a global
identity.

\subsection{An explicit nonzero certificate}

For $P(z)=z-1$, the coefficient bounds are $C=2,D=1$. Writing
$u=(t\sin t)^2$, the inner function is
\[
 A_P(t)=4\left((u-1)^2+
     16(1+u^2)^2\sin^2(\pi u)\right)-1.
\]
At the exact rational input
\[
 q=\frac{111415714087193}{100000000000000},
\]
the implementation produces the rigorous rational enclosure
\begin{equation}\label{eq:example-certificate}
 2-\frac{65}{2^{83}}\ \le E_P(q)\le2.
\end{equation}
No floating-point value is used in this certificate. The decimal-looking
input is parsed as an exact rational number; its proximity to a solution
of $t\sin t=1$ is a motivation for the example, not the justification
of the enclosure.

For $P(z)=z+1$, the only integer root is negative, so it has no natural
root. Because the construction uses $u=(t\sin t)^2\ge0$, it correctly
produces an identity. Indeed $P(u)\ge1$ immediately gives $A_P(t)\ge3$
and hence $E_P(t)=0$. This example checks why the witness domain must not
be silently changed from $\N$ to $\Z$.

For $P(z)=z^2+1$, the equation has no natural root and the same theorem
gives an identity. For the zero polynomial, the equation has natural
solutions, and the constructed expression is nonzero: at $t=0$ its value
is exactly $2$. These examples also make the direction of the reduction
clear: \emph{solvability corresponds to nonidentity}.

\subsection{Verification performed}

The implementation checks were run successfully: \textbf{14 checks passed}.
They cover rational interval operations, approximation widths, expression
construction, the example certificates, and a regression for evaluator
reuse with temporary expression nodes. An independent review found a
cache-identity lifetime bug in an earlier implementation; it was corrected
by retaining the expression object in every cached entry. The regression
reuses a single evaluator for 100 temporary rational-expression roots.
The corrected implementation was rerun independently.

These checks establish behavior of the supplied code and the indicated
individual certificates. They are not proofs of the global density
theorem, MRDP, or undecidability. The mathematical proof of those claims
is the argument in the report, with the explicitly named established
MRDP dependency. The new Lean source was not compiled. The LaTeX report
was compiled and its rendered pages inspected before delivery; the
build and verification records are included with the archive.

\subsection{Archive contents and commands}

\begin{center}
\begin{tabularx}{\linewidth}{@{}p{0.34\linewidth}X@{}}
\toprule
Path & Contents\\
\midrule
\code{report.tex}, \code{sections/} & Complete LaTeX source.\\
\code{report.pdf} & Compiled report.\\
\code{lean/} & Syntax, constructive logical core, and explicit object-formalization plan.\\
\code{code/} & Exact-rational compiler/evaluator, four examples, certificates, and implementation checks.\\
\code{evidence/} & Repository audit, source provenance, code review, and verification/build receipts.\\
\code{README.txt}, \code{build.sh} & Scope and reproduction instructions.\\
\bottomrule
\end{tabularx}
\end{center}

From the archive root, the numerical artifacts can be reproduced by
\begin{lstlisting}
python3 code/test_certificates.py
python3 code/richardson.py demo --output code/demo_results.json
python3 code/richardson.py evaluate \
  code/examples/linear_root.json --x 1.11415714087193 --bits 48
\end{lstlisting}
The PDF can be rebuilt with a working TeX Live installation using
\code{bash build.sh}. The Lean files are intended to be checked individually
with Lean 4 before beginning the larger ProveIt integration; their included
\code{\#print axioms} commands make their dependency inspection explicit.

No repository was modified or published as part of this work. The source
snapshot is used for analysis and an integration plan; the archive is a
standalone research deliverable.
